\scanpage{34}
à droite sur $k=\operatorname{End}(A)$, et ladite
antiéquivalence n'est autre que la
dualité de tels modules.

Revenons \deleted{\unsure{notre}}\inserted{à la} condition où
tt objet de $C$ est une $\varinjlim$ \unsure{dans|de} $C$
de type $\mathcal T'$ d'objets de $\mathcal A$. Alors
\unsure{tte|pr|toute|si} catégorie $C'$ où ces limites existent,
si $\underline{\operatorname{Hom}}_{\mathcal T'}(C,C')$ est la sous-catégorie
pleine de $\underline{\operatorname{Hom}}(C',C')$ formée des
foncteurs qui commutent à ces
$\varinjlim$, on trouve que
\[
 \underline{\operatorname{Hom}}_{\mathcal T'}(C,C')\longrightarrow\underline{\operatorname{Hom}}(\mathcal A,C')
\]
% slanted
\margin{\uline{à vérifier}}
est pl. fidèle \deleted{(et même une équival.
de cat. si \illegible{2} $C$ est libre}

Supposons $C$ additive, et $\mathcal T'$ l'ens
des cat. discrètes finies -- donc
tt objet de $C$ est somme d'objets
de $\mathcal A$, alors on trouve, si $C'$ additive
\[
 \underline{\operatorname{Hom}}_{\mathrm{add}}(C,C')\simeq
 \underline{\operatorname{Hom}}_{1/2\,\mathrm{add}}(\mathcal A,C')
\]
\unsure{Se rappelant|Si, rappelons|Si on applique} de l'inclusion
\[
 C\lhook\joinrel\xrightarrow{\mathrm{pl.\ fid.}}
 \deleted{\widehat{\mathcal A}_{\mathrm{ab}}}
 \underline{\operatorname{Hom}}_{1/2\,\mathrm{add}}(\mathcal A^\circ,\mathrm{Ab})
\]
et désignant par $D\subset\deleted{\widehat{\mathcal A}_{\mathrm{ab}}}\unsure{''}$ l'image
essentielle, on trouve alors
% archive sheet number: 33
